\section{Um teorema de existência para os feixes algébricos coerentes}
\label{section:III.5}


\subsection{Enunciado do teorema}
\label{subsection:III.5.1}

\begin{env}[5.1.1]
\label{III.5.1.1}
Seja $A$ um anel noetheriano ádico e seja $\mathfrak{J}$ um ideal de
definição de $A$, de modo que $A$ seja separado e completo para a
topologia $\mathfrak{J}$-ádica. Se $Y=\Spec(A)$, o pré-esquema formal
afim $\Spf(A)$ identifica-se com o completamento
$\widehat{Y}$ de $Y$ ao longo do subconjunto fechado
$Y'=\mathrm{V}(\mathfrak{J})$ \sref[I]{I.10.10.1}. Seja $X$ um
$Y$-pré-esquema ordinário de tipo finito, de morfismo estrutural
$f:X\to Y$. Denotemos por $\widehat{X}$ o completamento de $X$ ao longo
do subconjunto fechado $X'=f^{-1}(Y')$, ou, equivalentemente, o
$\widehat{Y}$-pré-esquema formal $X\times_Y\widehat{Y}$, e por
$\widehat{f}:\widehat{X}\to\widehat{Y}$ a extensão de $f$ às
completamentos. Finalmente, para todo $\sh{O}_X$-módulo coerente
$\sh{F}$, denotemos por $\widehat{\sh{F}}$ seu completamento
$\sh{F}_{/X'}$, que é um
$\sh{O}_{\widehat{X}}$-módulo coerente.
\end{env}

\begin{proposition}[5.1.2]
\label{III.5.1.2}
Com as hipóteses e a notação de \sref{III.5.1.1}, seja $\sh{F}$
um $\sh{O}_X$-módulo coerente cujo suporte é próprio sobre $Y$
\sref[II]{II.5.4.10}. Os homomorfismos canônicos
\sref{III.4.1.4}
\[
  \rho_i:
  \HH^i(X,\sh{F})
  \longrightarrow
  \HH^i(\widehat{X},\widehat{\sh{F}})
\]
são isomorfismos.
\end{proposition}

\begin{proof}
O $A$-módulo $\HH^i(X,\sh{F})$ é de tipo finito
\sref{III.3.2.4} e, portanto, identifica-se com seu completamento
separado para a topologia $\mathfrak{J}$-pré-ádica
\sref[0]{0.7.3.6}. A proposição é, por conseguinte, um caso particular
de \sref{III.4.1.10}.
\end{proof}

Recordemos que os isomorfismos canônicos $\rho_i$ comutam com os
homomorfismos de conexão associados a toda sucessão exata
\[
  0\longrightarrow\sh{F}'
  \longrightarrow\sh{F}
  \longrightarrow\sh{F}''
  \longrightarrow0
\]
de $\sh{O}_X$-módulos coerentes
\sref[0]{0.12.1.6} \sref[I]{I.10.8.9}.

\begin{corollary}[5.1.3]
\label{III.5.1.3}
Sejam $\sh{F}$ e $\sh{G}$ $\sh{O}_X$-módulos coerentes tais que
a interseção de seus suportes seja própria sobre $Y$. O
homomorfismo canônico
\[
  \Hom_{\sh{O}_X}(\sh{F},\sh{G})
  \longrightarrow
  \Hom_{\sh{O}_{\widehat{X}}}
  (\widehat{\sh{F}},\widehat{\sh{G}})
  \tag{5.1.3.1}
  \label{III.5.1.3.1.eq}
\]
que envia um homomorfismo $u:\sh{F}\to\sh{G}$ a seu completamento
$\widehat{u}:\widehat{\sh{F}}\to\widehat{\sh{G}}$, é um
isomorfismo. Além disso, quando $f$ é fechado,
$\widehat{u}$ é injetivo (respectivamente, sobrejetivo) se e somente
se $u$ é injetivo (respectivamente, sobrejetivo).
\end{corollary}

\begin{proof}
A primeira afirmação é um caso particular de \sref{III.4.5.3}; também
decorre do fato de que a fonte de \eqref{III.5.1.3.1.eq} é
um $A$-módulo de tipo finito e, portanto, identifica-se com seu
completamento separado. Para a segunda afirmação, observemos que, por
\sref[I]{I.10.8.14}, $\widehat{u}$ é injetivo (respectivamente,
sobrejetivo) se e somente se existe uma vizinhança de $X'$ em
$X$ sobre a qual $u$ é injetivo (respectivamente, sobrejetivo). A
conclusão deduz-se então do lema seguinte.
\end{proof}

\oldpage[III]{150}

\begin{lemma}[5.1.3.1]
\label{III.5.1.3.1}
Sob as hipóteses de \sref{III.5.1.1}, suponhamos além disso que
$f$ é fechado. Toda vizinhança de $X'$ em $X$ é então igual
a $X$.
\end{lemma}

\begin{proof}
Reduzimo-nos primeiro ao caso $f(X)=Y$. Com efeito, por hipótese,
$f(X)$ é um subconjunto fechado $Z$ de $Y$. Podemos substituir $f$ por
$f_{\mathrm{red}}$ \sref[I]{I.6.3.4} e supor assim que $X$
e $Y$ são reduzidos; podemos então substituir $Y$ pelo subpré-esquema
fechado reduzido cujo espaço subjacente é $Z$
\sref[I]{I.5.2.2}, pois todo ideal de $A$ é fechado e todo
anel quociente de $A$ é, portanto, ádico e noetheriano.

Tem-se então $f(X')=Y'$. Se $V$ é uma vizinhança aberta de $X'$
em $X$, o conjunto $f(X-V)$ é fechado em $Y$ e não intersecta $Y'$.
Isto é impossível, a menos que $X-V$ seja vazio, pois $\mathfrak{J}$
está contido no radical de $A$
\sref[I]{I.1.1.15} \sref[0]{0.7.1.10}. Portanto, $V=X$.
\end{proof}

Se nos restringimos aos $\sh{O}_X$-módulos coerentes cujo
suporte é próprio sobre $Y$, \sref{III.5.1.3} afirma, em linguagem
categórica, que o funtor
\[
  \sh{F}\longmapsto\widehat{\sh{F}}
\]
é plenamente fiel da categoria de tais $\sh{O}_X$-módulos para
a categoria dos $\sh{O}_{\widehat{X}}$-módulos coerentes. Fornece, assim,
uma equivalência da primeira categoria com uma subcategoria plena
da segunda \sref[0]{0.8.1.6}. O teorema de existência
mostrará que, quando $X$ é próprio sobre $Y$, esta subcategoria é em
realidade, a categoria de todos os
$\sh{O}_{\widehat{X}}$-módulos coerentes. Mais precisamente:

\begin{theorem}[5.1.4]
\label{III.5.1.4}
Seja $A$ um anel noetheriano ádico, seja $Y=\Spec(A)$, seja
$\mathfrak{J}$ um ideal de definição de $A$, e ponhamos
$Y'=\mathrm{V}(\mathfrak{J})$. Seja $f:X\to Y$ um morfismo separado
de tipo finito e ponhamos $X'=f^{-1}(Y')$. Sejam
\[
  \widehat{Y}=Y_{/Y'}=\Spf(A),\qquad
  \widehat{X}=X_{/X'}
\]
os completamentos de $Y$ e $X$ ao longo de $Y'$ e $X'$, e seja
$\widehat{f}:\widehat{X}\to\widehat{Y}$ a extensão de $f$ aos
completamentos. Então o funtor
\[
  \sh{F}\longmapsto\sh{F}_{/X'}=\widehat{\sh{F}}
\]
é uma equivalência da categoria dos
$\sh{O}_X$-módulos coerentes cujo suporte é próprio sobre $\Spec(A)$ com a
categoria dos $\sh{O}_{\widehat{X}}$-módulos coerentes cujo suporte
é próprio sobre $\Spf(A)$.
\end{theorem}

Em outras palavras, tendo em conta \sref{III.5.1.3}:

\begin{corollary}[5.1.5]
\label{III.5.1.5}
Um $\sh{O}_{\widehat{X}}$-módulo é isomorfo ao completamento de
um $\sh{O}_X$-módulo coerente cujo suporte é próprio sobre
$\Spec(A)$ se e somente se é coerente e seu suporte é próprio
sobre $\Spf(A)$.
\end{corollary}

O caso mais importante é o seguinte:

\begin{corollary}[5.1.6]
\label{III.5.1.6}
Suponhamos que $X$ é próprio sobre $Y=\Spec(A)$. Então o funtor
$\sh{F}\mapsto\widehat{\sh{F}}$ é uma equivalência entre a
categoria dos $\sh{O}_X$-módulos coerentes e a categoria dos
$\sh{O}_{\widehat{X}}$-módulos coerentes.
\end{corollary}

\begin{env}[5.1.7]
\label{III.5.1.7}
\emph{Escólio.}
Tendo em conta a caracterização dos feixes coerentes sobre
pré-esquemas formais \sref[I]{I.10.11.3}, vemos que, sob as
condições de \sref{III.5.1.1}, dar um
$\sh{O}_X$-módulo coerente cujo suporte é próprio sobre $\Spec(A)$
equivale---se
\[
  Y_n=\Spec(A/\mathfrak{J}^{n+1}),\qquad
  X_n=X\times_Y Y_n
\]
---a dar um sistema projetivo $(\sh{F}_n)$ de
$\sh{O}_{X_n}$-módulos coerentes tal que, para $m\leq n$,
\[
  \sh{F}_m
  =
  \sh{F}_n\otimes_{\sh{O}_{Y_n}}\sh{O}_{Y_m},
  \qquad\text{ou, equivalentemente}\qquad
  \sh{F}_m
  =
  \sh{F}_n/\mathfrak{J}^{m+1}\sh{F}_n,
\]
e tal que o suporte de $\sh{F}_0$ seja um subconjunto de $X_0$
próprio sobre $Y_0$. Por \sref[I]{I.10.11.4}, os homomorfismos de
$\sh{O}_X$-módulos coerentes são interpretados de modo análogo como
homomorfismos de sistemas projetivos de
$\sh{O}_{X_n}$-módulos coerentes.

\oldpage[III]{151}

Em todas as aplicações conhecidas, $A$ é de fato um anel
noetheriano ádico local. Assim, os $Y_n$ são espectros de anéis
artinianos locais, e os resultados desta subseção e das
anteriores reduzem, em grande medida, a geometria algébrica sobre um
anel noetheriano ádico local à geometria algébrica sobre anéis
artinianos locais.
\end{env}

\begin{corollary}[5.1.8]
\label{III.5.1.8}
Sob as hipóteses de \sref{III.5.1.4}, a aplicação
\[
  Z\longmapsto\widehat{Z}=Z_{/(Z\cap X')}
\]
é uma bijeção do conjunto dos subpré-esquemas fechados $Z$ de $X$
que são próprios sobre $Y$ com o conjunto dos subpré-esquemas
formais fechados de $\widehat{X}$ que são próprios sobre
$\widehat{Y}$.
\end{corollary}

\begin{proof}
Um subpré-esquema formal fechado de $\widehat{X}$ é da forma
\[
  \bigl(T,
    (\sh{O}_{\widehat{X}}/\sh{A})|_T\bigr),
\]
onde $\sh{A}$ é um ideal coerente de
$\sh{O}_{\widehat{X}}$ \sref[I]{I.10.14.2}. Se $T$ é próprio sobre
$\widehat{Y}$, \sref{III.5.1.4} mostra que
$\sh{O}_{\widehat{X}}/\sh{A}$ é isomorfo a um
$\sh{O}_{\widehat{X}}$-módulo da forma
$\widehat{\sh{F}}$, onde $\sh{F}$ é um
$\sh{O}_X$-módulo coerente cujo suporte é próprio sobre $Y$. Além disso,
\sref{III.5.1.3} mostra que o homomorfismo canônico
\[
  \sh{O}_{\widehat{X}}
  \longrightarrow
  \sh{O}_{\widehat{X}}/\sh{A}
\]
é da forma $\widehat{u}$, onde
$u:\sh{O}_X\to\sh{F}$ é um homomorfismo sobrejetivo de
$\sh{O}_X$-módulos. Por conseguinte,
$\sh{F}=\sh{O}_X/\sh{N}$ para um ideal coerente $\sh{N}$ de
$\sh{O}_X$, e
$\sh{A}=\widehat{\sh{N}}$ \sref[I]{I.10.8.8}. A conclusão se
deduz agora de \sref[I]{I.10.14.7}.
\end{proof}


\subsection{Demonstração do teorema de existência: caso projetivo e quase-projetivo}
\label{subsection:III.5.2}

\begin{env}[5.2.1]
\label{III.5.2.1}
Sob as condições de \sref{III.5.1.4}, diremos
provisoriamente que um $\sh{O}_{\widehat{X}}$-módulo coerente é
\emph{algebrizável} se é isomorfo ao completamento
$\widehat{\sh{F}}$ de um $\sh{O}_X$-módulo coerente $\sh{F}$ cujo
suporte é próprio sobre $Y$.
\end{env}

\begin{lemma}[5.2.2]
\label{III.5.2.2}
Sejam $\sh{F}'$ e $\sh{G}'$ $\sh{O}_{\widehat{X}}$-módulos
algebrizáveis. Para todo homomorfismo
$u:\sh{F}'\to\sh{G}'$, os módulos
$\Ker(u)$, $\Im(u)$ e $\Coker(u)$ são algebrizáveis.
\end{lemma}

\begin{proof}
Escrevamos
\[
  \sh{F}'=\widehat{\sh{F}},\qquad
  \sh{G}'=\widehat{\sh{G}},
\]
onde $\sh{F}$ e $\sh{G}$ são $\sh{O}_X$-módulos coerentes cujos
suportes são próprios sobre $Y$. Por \sref{III.5.1.3},
$u=\widehat{v}$ para um homomorfismo
$v:\sh{F}\to\sh{G}$. Como o completamento é exato,
$\Ker(\widehat{v})\simeq\widehat{\Ker(v)}$; e como o suporte
de $\Ker(v)$ está contido no de $\sh{F}$, deduz-se que
$\Ker(u)$ é algebrizável. Os argumentos para $\Im(u)$ e
$\Coker(u)$ são os mesmos.
\end{proof}

\begin{proposition}[5.2.3]
\label{III.5.2.3}
Seja $A$ um anel noetheriano ádico, seja $\mathfrak{J}$ um ideal de
definição de $A$, seja $\mathfrak{Y}=\Spf(A)$, e seja
$f:\mathfrak{X}\to\mathfrak{Y}$ um morfismo próprio de pré-esquemas
formais. Ponhamos
\[
  Y_k=\Spec(A/\mathfrak{J}^{k+1}),\qquad
  X_k=\mathfrak{X}\times_{\mathfrak{Y}}Y_k,
\]
e, para todo $\sh{O}_{\mathfrak{X}}$-módulo $\sh{F}$, ponhamos
\[
  \sh{F}_k
  =
  \sh{F}\otimes_{\sh{O}_{\mathfrak{X}}}\sh{O}_{X_k}
  =
  \sh{F}/\mathfrak{J}^{k+1}\sh{F}.
\]
Seja $\sh{L}$ um $\sh{O}_{\mathfrak{X}}$-módulo invertível, e
suponhamos que
$\sh{L}_0=\sh{L}/\mathfrak{J}\sh{L}$ é um
$\sh{O}_{X_0}$-módulo amplo. Para todo
$\sh{O}_{\mathfrak{X}}$-módulo $\sh{F}$ e todo inteiro $n$, ponhamos
\[
  \sh{F}(n)=\sh{F}\otimes\sh{L}^{\otimes n}.
\]
Então, para todo $\sh{O}_{\mathfrak{X}}$-módulo coerente $\sh{F}$,
existe um inteiro $n_0$ tal que, para todo $n\geq n_0$:
\begin{enumerate}
  \item[\rm{(i)}] o homomorfismo canônico
    \[
      \HH^0(\mathfrak{X},\sh{F}(n))
      \longrightarrow
      \HH^0(\mathfrak{X},\sh{F}_k(n))
    \]
    é sobrejetivo para todo $k\geq0$;
  \item[\rm{(ii)}]
    $\HH^q(\mathfrak{X},\sh{F}(n))=0$ para todo $q>0$.
\end{enumerate}
\end{proposition}

\begin{proof}
Os espaços subjacentes a $\mathfrak{X}$ e $X_0$ são os mesmos.
Os feixes
\[
  \sh{M}_k
  =
  \mathfrak{J}^k\sh{F}/\mathfrak{J}^{k+1}\sh{F},
\]
por serem anulados por $\mathfrak{J}$, podem ser considerados
$\sh{O}_{X_0}$-módulos coerentes \sref[0]{0.5.3.10}. Além disso, se
\[
  \sh{M}_k(n)
  =
  \sh{M}_k\otimes_{\sh{O}_{X_0}}\sh{L}_0^{\otimes n},
\]
então
\[
  \sh{M}_k(n)
  =
  \mathfrak{J}^k\sh{F}(n)/
  \mathfrak{J}^{k+1}\sh{F}(n).
\]
Como $\sh{L}_0$ é amplo para $f_0:X_0\to Y_0$ e $f_0$ é
próprio,
\oldpage[III]{152}
$f_0$ é projetivo \sref[II]{II.5.5.4}. Seja
\[
  S=\bigoplus_{k\geq0}
  \mathfrak{J}^k/\mathfrak{J}^{k+1}
\]
a $A$-álgebra graduada associada à
filtração $\mathfrak{J}$-ádica de $A$. É de tipo finito
porque $A$ é noetheriano. Ponhamos
\[
  \sh{S}'=f_0^*(\widetilde{S}),\qquad
  \sh{M}=\bigoplus_{k\geq0}\sh{M}_k.
\]
Então $\sh{S}'$ é uma $\sh{O}_{X_0}$-álgebra quase-coerente de
tipo finito, e $\sh{M}$ é um $\sh{S}'$-módulo graduado
quase-coerente de tipo finito, pois $\sh{F}_0$ é coerente e
gera $\sh{M}$ como $\sh{S}'$-módulo. O Teorema
\sref{III.2.4.1},~\emph{(ii)}, fornece, portanto, um inteiro $n_0$
tal que, para $n\geq n_0$ e todo $k$,
\[
  \HH^q(X_0,\sh{M}_k(n))=0
  \qquad(q>0).
  \tag{5.2.3.1}
  \label{III.5.2.3.1}
\]
Considerado agora como $\sh{O}_{\mathfrak{X}}$-módulo,
$\sh{M}_k(n)$ satisfaz, por conseguinte, também
\[
  \HH^q(\mathfrak{X},\sh{M}_k(n))=0
  \qquad(q>0,\ n\geq n_0).
\]

Apliquemos a sucessão exata de cohomologia a
\[
  0\longrightarrow
  \frac{\mathfrak{J}^k\sh{F}(n)}
       {\mathfrak{J}^{k+1}\sh{F}(n)}
  \longrightarrow
  \frac{\mathfrak{J}^h\sh{F}(n)}
       {\mathfrak{J}^{k+1}\sh{F}(n)}
  \longrightarrow
  \frac{\mathfrak{J}^h\sh{F}(n)}
       {\mathfrak{J}^{k}\sh{F}(n)}
  \longrightarrow0.
\]
Para $0\leq h<k$, $n\geq n_0$ e $q>0$, uma indução sobre $k-h$
dá
\[
  \HH^q\left(
    \mathfrak{X},
    \frac{\mathfrak{J}^h\sh{F}(n)}
         {\mathfrak{J}^{k}\sh{F}(n)}
  \right)=0,
  \tag{5.2.3.2}
  \label{III.5.2.3.2}
\]
e, em particular, para $h=0$,
\[
  \HH^q(\mathfrak{X},\sh{F}_k(n))=0
  \qquad(n\geq n_0,\ k\geq0,\ q>0).
  \tag{5.2.3.3}
  \label{III.5.2.3.3}
\]
Outra parte da sucessão exata de cohomologia, com $h=0$, dá
\[
  \HH^0(\mathfrak{X},\sh{F}_{k+1}(n))
  \longrightarrow
  \HH^0(\mathfrak{X},\sh{F}_k(n))
  \longrightarrow
  \HH^1\left(
    \mathfrak{X},
    \frac{\mathfrak{J}^k\sh{F}(n)}
         {\mathfrak{J}^{k+1}\sh{F}(n)}
  \right)
  =0.
  \tag{5.2.3.4}
  \label{III.5.2.3.4}
\]
Deduz-se que, para $h\leq k$, a aplicação canônica
\[
  \HH^0(\mathfrak{X},\sh{F}_k(n))
  \longrightarrow
  \HH^0(\mathfrak{X},\sh{F}_h(n))
  \tag{5.2.3.5}
  \label{III.5.2.3.5}
\]
é sobrejetiva. Assim, para todo $q$, o sistema projetivo
\[
  \bigl(\HH^q(\mathfrak{X},\sh{F}_k(n))\bigr)_{k\geq0}
\]
satisfaz a condição (ML) quando $n\geq n_0$.

Além disso, todo aberto formal afim $U$ de $\mathfrak{X}$ é um
aberto afim de cada $X_k$ \sref[I]{I.10.5.2}. Portanto,
$\HH^q(U,\sh{F}_k(n))=0$ para $q>0$ \sref[I]{I.1.3.1}, e
\[
  \HH^0(U,\sh{F}_k(n))
  \longrightarrow
  \HH^0(U,\sh{F}_h(n))
\]
é sobrejetivo para $h\leq k$ \sref[I]{I.1.3.9}. Portanto, são satisfeitas
as hipóteses de \sref[0]{0.13.3.1}. Para $n\geq n_0$
obtemos:
\begin{enumerate}
  \item[$1^\circ$] para todo $q>0$, a aplicação
    \[
      \HH^q(\mathfrak{X},\sh{F}(n))
      \longrightarrow
      \varprojlim_k
      \HH^q(\mathfrak{X},\sh{F}_k(n))
    \]
    é bijetiva; portanto, por \eqref{III.5.2.3.3},
    $\HH^q(\mathfrak{X},\sh{F}(n))=0$;
  \item[$2^\circ$] o homomorfismo
    \[
      \HH^0(\mathfrak{X},\sh{F}(n))
      \longrightarrow
      \varprojlim_k
      \HH^0(\mathfrak{X},\sh{F}_k(n))
    \]
    é bijetivo. Como os homomorfismos
    \eqref{III.5.2.3.5} são sobrejetivos, também o é cada homomorfismo
    \[
      \varprojlim_k
      \HH^0(\mathfrak{X},\sh{F}_k(n))
      \longrightarrow
      \HH^0(\mathfrak{X},\sh{F}_h(n)).
    \]
\end{enumerate}
Isto demonstra as duas afirmações.
\end{proof}

\begin{corollary}[5.2.4]
\label{III.5.2.4}
Sob as hipóteses de \sref{III.5.2.3}, para todo
$\sh{O}_{\mathfrak{X}}$-módulo coerente $\sh{F}$ existe um inteiro $N$
tal que, para $n\geq N$, $\sh{F}(n)$ está gerado por suas seções
sobre
\oldpage[III]{153}
$\mathfrak{X}$. Equivalentemente, $\sh{F}$ é isomorfo a um quociente
de um $\sh{O}_{\mathfrak{X}}$-módulo da forma
$(\sh{O}_{\mathfrak{X}}(-n))^k$.
\end{corollary}

\begin{proof}
Como $X_0$ é noetheriano, a hipótese sobre $\sh{L}_0$ e
\sref[II]{II.4.5.5} fornecem um inteiro $n_0$ tal que
$\sh{F}_0(n)$ está gerado por suas seções sobre $\mathfrak{X}$ para
$n\geq n_0$. Podemos tomar $n_0$ suficientemente grande para que
\[
  \Gamma(\mathfrak{X},\sh{F}(n))
  \longrightarrow
  \Gamma(\mathfrak{X},\sh{F}_0(n))
\]
seja também sobrejetivo para $n\geq n_0$ \sref{III.5.2.3}. Existem,
portanto, um número finito de seções
$s_i\in\Gamma(\mathfrak{X},\sh{F}(n))$ cujas imagens geram
$\sh{F}_0(n)$ \sref[0]{0.5.2.3}. Como $\mathfrak{J}$ está
contido no ideal maximal do anel local em todo ponto
de $\mathfrak{X}$, o lema de Nakayama aplicado a esses anéis locais
mostra que os $s_i$ geram $\sh{F}(n)$
\sref[0]{0.5.1.1}.
\end{proof}

\begin{env}[5.2.5]
\label{III.5.2.5}
\emph{Demonstração do teorema de existência: caso projetivo.}
Com a notação de \sref{III.5.1.4}, suponhamos que $f$ é
projetivo, de modo que existe um $\sh{O}_X$-módulo amplo $\sh{L}$
\sref[I]{I.5.5.4}. Por definição,
\[
  X_n=\widehat{X}\times_{\widehat{Y}}Y_n
\]
é o subpré-esquema fechado
$X\times_Y Y_n=f^{-1}(Y_n)$ de $X$. Se
\[
  \sh{L}'=\sh{L}\otimes_{\sh{O}_X}\sh{O}_{\widehat{X}}
\]
é o completamento de $\sh{L}$, então
$\sh{L}'_0=\sh{L}/\mathfrak{J}\sh{L}$, considerado como
$\sh{O}_{X_0}$-módulo; este módulo é amplo
\sref[II]{II.4.6.13}[(i~\emph{bis})].

Podemos, portanto, aplicar \sref{III.5.2.4} a $\sh{L}'$ e a todo
$\sh{O}_{\widehat{X}}$-módulo coerente $\sh{F}$. Deduz-se que
$\sh{F}$ é isomorfo a um quociente de
\[
  \sh{G}=
  \bigl(\sh{L}'^{\otimes(-n)}\bigr)^k
\]
para inteiros adequados $n>0$ e $k>0$. O módulo $\sh{G}$ é o
completamento de $(\sh{L}^{\otimes(-n)})^k$
\sref[I]{I.10.8.10} e, portanto, é algebrizável. Seja
$u:\sh{G}\to\sh{F}$ a sobrejeção canônica, e ponhamos
$\sh{H}=\Ker(u)$; este é um
$\sh{O}_{\widehat{X}}$-módulo coerente \sref[0]{0.5.3.4}. Aplicando o mesmo
argumento a $\sh{H}$, obtemos um
$\sh{O}_{\widehat{X}}$-módulo algebrizável $\sh{K}$ e um homomorfismo
$v:\sh{K}\to\sh{G}$ tal que $\sh{H}=\Im(v)$. Assim,
$\sh{F}=\Coker(v)$, e $\sh{F}$ é algebrizável por
\sref{III.5.2.2}.
\end{env}

\begin{env}[5.2.6]
\label{III.5.2.6}
\emph{Demonstração do teorema de existência: caso quase-projetivo.}
Conservemos a notação de \sref{III.5.1.4} e suponhamos agora que $f$
é quase-projetivo. Existe um morfismo projetivo
$g:Z\to Y$ tal que $X$ identifica-se com o $Y$-pré-esquema
induzido sobre um subconjunto aberto de $Z$ \sref[II]{II.5.3.2}. Se
$Z'=g^{-1}(Y')$, então $X'=X\cap Z'$. Por conseguinte, o completamento
$\widehat{X}=X_{/X'}$ identifica-se com o pré-esquema formal
induzido por $\widehat{Z}=Z_{/Z'}$ sobre o subconjunto aberto $X\cap Z'$ de
$\widehat{Z}$ \sref[I]{I.10.8.5}.

Seja $\sh{F}'$ um $\sh{O}_{\widehat{X}}$-módulo coerente cujo
suporte $T'$ seja próprio sobre $\widehat{Y}$. Por definição, existe
um subpré-esquema fechado de $X'$ de espaço subjacente
$T'\subset X'$ tal que a restrição $T'\to Y'$ de $f$ é
própria. Portanto, $T'$ é próprio sobre $Y$ e, consequentemente, fechado em
$Z'$ \sref[II]{II.5.4.10}. Deduz-se que $\sh{F}'$ é induzido
sobre $\widehat{X}$ pelo $\sh{O}_{\widehat{Z}}$-módulo $\sh{G}'$
obtido pela colagem de $\sh{F}'$, sobre o aberto $\widehat{X}$, com o
feixe zero sobre o aberto $\widehat{Z}-T'$; ambos os feixes coincidem sobre
sua interseção $\widehat{X}-T'$.

O módulo $\sh{G}'$ é coerente. Por \sref{III.5.2.5}, existe um
$\sh{O}_Z$-módulo coerente $\sh{G}$ tal que
$\sh{G}'=\widehat{\sh{G}}$. Seja $T$ o suporte de $\sh{G}$;
então $T'=T\cap Z'$ \sref[I]{I.10.8.12}. Se $h$ é a restrição
de $g$ ao subpré-esquema fechado reduzido de $Z$ cujo espaço subjacente
é $T$, então
\[
  T'=h^{-1}(Y')=T\cap g^{-1}(Y').
\]
Assim, $X\cap T$ é um subconjunto aberto de $T$ que contém $T'$.

\oldpage[III]{154}
Como $h$ é próprio \sref[II]{II.5.4.2}, e portanto fechado,
\sref{III.5.1.3.1} dá $X\cap T=T$. Assim, $T\subset X$; e
como $T$ é fechado em $Z$, é próprio sobre $Y$. Se $\sh{F}$ é
o feixe induzido sobre $X$ por $\sh{G}$, seu completamento
$\widehat{\sh{F}}$ é induzido sobre $\widehat{X}$ por
$\widehat{\sh{G}}$ \sref[I]{I.10.8.4} e, portanto, é igual a
$\sh{F}'$. Isto termina a demonstração.
\end{env}


\subsection{Demonstração do teorema de existência: caso geral}
\label{subsection:III.5.3}

\begin{lemma}[5.3.1]
\label{III.5.3.1}
Sob as condições de \sref{III.5.1.4}, seja
\[
  0\longrightarrow\sh{H}\longrightarrow\sh{F}
  \longrightarrow\sh{G}\longrightarrow0
\]
uma sucessão exata de $\sh{O}_{\widehat{X}}$-módulos coerentes.
Se $\sh{G}$ e $\sh{H}$ são algebrizáveis, então $\sh{F}$ é
algebrizável.
\end{lemma}

\begin{proof}
Suponhamos que
\[
  \sh{G}=\widehat{\sh{B}},
  \qquad
  \sh{H}=\widehat{\sh{C}},
\]
onde $\sh{B}$ e $\sh{C}$ são $\sh{O}_X$-módulos coerentes cujos
suportes são próprios sobre $Y$. O módulo $\sh{F}$ define canonicamente
um elemento do $A$-módulo
\[
  \Ext^1_{\sh{O}_{\widehat{X}}}
  (\widehat{X};\widehat{\sh{B}},\widehat{\sh{C}})
\]
\sref[0]{0.12.3.2}. As hipóteses implicam que este $A$-módulo é
canonicamente isomorfo a
\[
  \Ext^1_{\sh{O}_X}(X;\sh{B},\sh{C})
\]
\sref{III.4.5.2}. Existe, portanto, uma sucessão exata
\[
  0\longrightarrow\sh{C}\longrightarrow\sh{A}
  \longrightarrow\sh{B}\longrightarrow0
\]
de $\sh{O}_X$-módulos coerentes cuja classe canônica é enviada para a
classe correspondente a $\sh{F}$. Por definição, tendo em conta
\sref[I]{I.10.8.8}, (ii), isto significa que $\sh{F}$ é isomorfo a
$\widehat{\sh{A}}$. O lema segue, pois,
$\Supp(\sh{A})$ está contido em
$\Supp(\sh{B})\cup\Supp(\sh{C})$ e é, portanto, próprio sobre $Y$.
\end{proof}

\begin{corollary}[5.3.2]
\label{III.5.3.2}
Sob as condições de \sref{III.5.1.1}, seja
$u:\sh{F}\to\sh{G}$ um homomorfismo de
$\sh{O}_{\widehat{X}}$-módulos coerentes. Se $\sh{G}$, $\Ker(u)$ e
$\Coker(u)$ são algebrizáveis, então $\sh{F}$ é algebrizável.
\end{corollary}

\begin{proof}
O Lema \sref{III.5.2.2}, aplicado a
$\sh{G}\to\Coker(u)$, mostra que $\Im(u)$ é algebrizável. Basta então
aplicar o Lema \sref{III.5.3.1} à sucessão exata
\[
  0\longrightarrow\Ker(u)\longrightarrow\sh{F}
  \longrightarrow\Im(u)\longrightarrow0.
\]
\end{proof}

\begin{lemma}[5.3.3]
\label{III.5.3.3}
Sob as condições de \sref{III.5.1.1}, seja $h:Z\to Y$ um
morfismo de tipo finito, seja $\widehat{Z}$ o completamento de $Z$
ao longo de $Z'=h^{-1}(Y')$, seja $g:Z\to X$ um $Y$-morfismo próprio,
e seja $\widehat{g}:\widehat{Z}\to\widehat{X}$ sua extensão aos
completamentos. Para todo $\sh{O}_{\widehat{Z}}$-módulo
algebrizável $\sh{F}'$, o módulo
$\widehat{g}_*(\sh{F}')$ é um
$\sh{O}_{\widehat{X}}$-módulo algebrizável.
\end{lemma}

\begin{proof}
Se $\sh{F}$ é um $\sh{O}_Z$-módulo coerente tal que
$\sh{F}'=\widehat{\sh{F}}$, o primeiro teorema de comparação
\sref{III.4.1.5} identifica
$\widehat{g}_*(\widehat{\sh{F}})$ com o completamento de
$g_*(\sh{F})$.
\end{proof}

\begin{lemma}[5.3.4]
\label{III.5.3.4}
Seja $X$ um esquema ordinário noetheriano, seja $X'$ um subconjunto
fechado de $X$, seja $f:Z\to X$ um morfismo próprio, e ponhamos
\[
  Z'=f^{-1}(X'),\qquad
  \widehat{X}=X_{/X'},\qquad
  \widehat{Z}=Z_{/Z'}.
\]
Seja $\widehat{f}:\widehat{Z}\to\widehat{X}$ a extensão de $f$
aos completamentos. Seja $\sh{M}$ um ideal coerente de
$\sh{O}_X$ tal que, se
\[
  U=X-\Supp(\sh{O}_X/\sh{M}),
\]
a restrição $f^{-1}(U)\to U$ de $f$ é um isomorfismo. Então,
para todo $\sh{O}_{\widehat{X}}$-módulo coerente $\sh{F}$, existe
um inteiro $n>0$ tal que o núcleo e o conúcleo do homomorfismo
canônico
\[
  \rho_{\sh{F}}:
  \sh{F}\longrightarrow
  \widehat{f}_*\bigl(\widehat{f}^{\,*}(\sh{F})\bigr)
\]
\sref[0]{0.4.4.3} são anulados por
$\widehat{\sh{M}}^{\,n}$.
\end{lemma}

\begin{proof}
Podemos restringir-nos ao caso $X=\Spec(B)$, onde $B$ é um
anel noetheriano; então $X'=\mathrm{V}(\mathfrak{K})$ para um ideal
$\mathfrak{K}$ de $B$. Reduzimo-nos primeiro ao caso em que $B$ é
noetheriano ádico e $\mathfrak{K}$ é um ideal de definição.

Seja $B_1$ o completamento separado de $B$ para a
topologia $\mathfrak{K}$-pré-ádica.
\oldpage[III]{155}
Se $\mathfrak{K}_1=\mathfrak{K}B_1$, então $B_1$ é um anel
noetheriano ádico e $\mathfrak{K}_1$ é um ideal de definição.
Ponhamos $X_1=\Spec(B_1)$, e seja $h:X_1\to X$ o morfismo
correspondente ao homomorfismo canônico $B\to B_1$. Se
$X'_1=h^{-1}(X')$, então $X'_1=\mathrm{V}(\mathfrak{K}_1)$. Finalmente,
ponhamos
\[
  Z_1=Z\times_X X_1,\qquad f_1:Z_1\longrightarrow X_1.
\]
O morfismo $f_1$ é próprio \sref[II]{II.5.4.2}. Denotemos por
$\widehat{X}_1$ o completamento de $X_1$ ao longo de $X'_1$, por
\[
  \widehat{Z}_1=
  Z_1\times_{X_1}\widehat{X}_1
\]
o completamento de $Z_1$ ao longo de
$Z'_1=f_1^{-1}(X'_1)$, e por
$\widehat{f}_1$ a extensão de $f_1$ aos completamentos.

A extensão
$\widehat{h}:\widehat{X}_1\to\widehat{X}$ corresponde à
aplicação identidade de $B_1$ e é, portanto, um isomorfismo
\sref[I]{I.10.9.1}. O morfismo correspondente
$\widehat{Z}_1\to\widehat{Z}$ é também, por conseguinte, um isomorfismo,
e estes isomorfismos identificam $\widehat{f}_1$ com
$\widehat{f}$. Além disso,
$\sh{M}_1=h^*(\sh{M})$ é um ideal coerente de $\sh{O}_{X_1}$ e
\[
  \Supp(\sh{O}_{X_1}/\sh{M}_1)
  =h^{-1}\bigl(\Supp(\sh{O}_X/\sh{M})\bigr)
\]
\sref[I]{I.9.1.13}. Assim, se
$U_1=X_1-\Supp(\sh{O}_{X_1}/\sh{M}_1)$, então
$U_1=h^{-1}(U)$, e a restrição
$f_1^{-1}(U_1)\to U_1$ é um isomorfismo
\sref[I]{I.3.2.7}. Finalmente, $\widehat{\sh{M}}$ e
$\widehat{\sh{M}}_1$ são identificados por $\widehat{h}$
\sref[I]{I.10.9.5}. Todas as hipóteses do lema são, pois,
satisfeitas por $X_1$, $X'_1$, $f_1$ e $\sh{M}_1$; doravante
podemos supor que $B$ é noetheriano ádico e que $\mathfrak{K}$ é
um ideal de definição.

Tem-se então $\widehat{X}=\Spf(B)$ e
$\sh{F}=N^\Delta$, onde $N$ é um $B$-módulo de tipo finito. Assim,
$\sh{F}=\widehat{\sh{G}}$, onde $\sh{G}$ é o
$\sh{O}_X$-módulo coerente $\widetilde{N}$ \sref[I]{I.10.10.5}, e
\[
  \widehat{f}^{\,*}(\sh{F})
  =\widehat{f^*(\sh{G})}
\]
\sref[I]{I.10.9.5}. Pelo primeiro teorema de comparação
\sref{III.4.1.5},
\[
  \widehat{f}_*
  \bigl(\widehat{f^*(\sh{G})}\bigr)
  \simeq
  \widehat{f_*\bigl(f^*(\sh{G})\bigr)},
\]
e, por \sref{III.5.1.3}, o homomorfismo canônico
$\rho_{\sh{F}}$ é precisamente $\widehat{\rho}_{\sh{G}}$.

O núcleo $\sh{P}$ e o conúcleo $\sh{R}$ de
\[
  \rho_{\sh{G}}:
  \sh{G}\longrightarrow f_*\bigl(f^*(\sh{G})\bigr)
\]
são $\sh{O}_X$-módulos coerentes, e suas restrições a $U$ são
nulas. Existe, portanto, um inteiro $n>0$ tal que
\[
  \sh{M}^n\sh{P}=\sh{M}^n\sh{R}=0
\]
\sref[I]{I.9.3.4}. Deduz-se que
\[
  \widehat{\sh{M}}^{\,n}\widehat{\sh{P}}
  =
  \widehat{\sh{M}}^{\,n}\widehat{\sh{R}}
  =0
\]
\sref[I]{I.10.8.8} \sref[I]{I.10.8.10}, o que demonstra o lema.
\end{proof}

\begin{env}[5.3.5]
\label{III.5.3.5}
\emph{Fim da demonstração do teorema de existência.}
Conservemos as hipóteses de \sref{III.5.1.4}. Utilizemos a indução
noetheriana \sref[0]{0.2.2.2}, supondo conhecido o teorema para todo
subpré-esquema fechado $T$ de $X$ cujo espaço subjacente é distinto
de $X$; aqui $\widehat{T}$ denota o completamento de $T$ ao longo de
$T\cap X'$. Podemos supor que $X$ não é vazio.

Como $f$ é separado e de tipo finito, o lema de Chow
\sref[II]{II.5.6.1} fornece um $Y$-esquema $Z$ e um $Y$-morfismo
$g:Z\to X$ tais que o morfismo estrutural $h:Z\to Y$ é
quase-projetivo, $g$ é projetivo e sobrejetivo, e existe um
subconjunto aberto não vazio $U$ de $X$ no qual
$g^{-1}(U)\to U$ é um isomorfismo. Seja $\sh{M}$ um ideal coerente
de $\sh{O}_X$ que define um subpré-esquema fechado de espaço subjacente
$X-U$ \sref[I]{I.5.2.2}, e seja $\sh{F}$ um
$\sh{O}_{\widehat{X}}$-módulo coerente cujo suporte $E$ é próprio sobre $Y$.
Denotemos por $\widehat{Z}$ o completamento de $Z$ ao longo de
$h^{-1}(Y')$, e por
$\widehat{g}:\widehat{Z}\to\widehat{X}$ a extensão de $g$ aos
completamentos.

O módulo $\widehat{g}^{\,*}(\sh{F})$ é um
$\sh{O}_{\widehat{Z}}$-módulo coerente cujo suporte está contido em
$g^{-1}(E)$ e é, portanto, próprio sobre $Y$, pois $g$ é
projetivo e, portanto, próprio \sref[II]{II.5.4.6}. Como $h$ é
quase-projetivo, $\widehat{g}^{\,*}(\sh{F})$ é algebrizável por
\sref{III.5.2.6}.
\oldpage[III]{156}
Deduz-se de \sref{III.5.3.3}, dado que $g$ é próprio, que
\[
  \widehat{g}_*
  \bigl(\widehat{g}^{\,*}(\sh{F})\bigr)
\]
é um $\sh{O}_{\widehat{X}}$-módulo algebrizável.

Apliquemos \sref{III.5.3.4} a $\sh{F}$ e $g$. O núcleo $\sh{P}$
e o conúcleo $\sh{R}$ de
\[
  \rho_{\sh{F}}:
  \sh{F}\longrightarrow
  \widehat{g}_*\bigl(\widehat{g}^{\,*}(\sh{F})\bigr)
\]
são anulados por uma potência $\widehat{\sh{M}}^{\,n}$. Seja $T$
o subpré-esquema fechado de $X$ definido por $\sh{M}^n$, cujo
espaço subjacente é $X-U$, e seja $j:T\to X$ a imersão
canônica. Sua extensão
$\widehat{j}:\widehat{T}\to\widehat{X}$ é a imersão canônica
\sref[I]{I.10.14.7}, e podemos escrever
\[
  \sh{P}=
  \widehat{j}_*\bigl(\widehat{j}^{\,*}(\sh{P})\bigr),
  \qquad
  \sh{R}=
  \widehat{j}_*\bigl(\widehat{j}^{\,*}(\sh{R})\bigr).
\]
Como $U$ não é vazio, a hipótese de indução mostra que
$\widehat{j}^{\,*}(\sh{P})$ e
$\widehat{j}^{\,*}(\sh{R})$ são $\sh{O}_{\widehat{T}}$-módulos
algebrizáveis. O Lema \sref{III.5.3.3} mostra então
que $\sh{P}$ e $\sh{R}$ são algebrizáveis, e
\sref{III.5.3.2} demonstra finalmente que $\sh{F}$ é algebrizável.
\end{env}


\subsection{Aplicação: comparação entre morfismos de pré-esquemas ordinários e morfismos de pré-esquemas formais. Pré-esquemas formais algebrizáveis}
\label{subsection:III.5.4}

\begin{theorem}[5.4.1]
\label{III.5.4.1}
Seja $A$ um anel noetheriano ádico, seja $\mathfrak{J}$ um ideal
de definição de $A$, e ponhamos
\[
  S=\Spec(A),\qquad S'=\mathrm{V}(\mathfrak{J}).
\]
Seja $u:X\to S$ um morfismo próprio e seja $v:Y\to S$ um
morfismo separado de tipo finito. Sejam
$\widehat{S}$, $\widehat{X}$ e $\widehat{Y}$ os completamentos
de $S$, $X$ e $Y$ ao longo de $S'$, $u^{-1}(S')$ e $v^{-1}(S')$,
respectivamente. Para todo $S$-morfismo $f:X\to Y$, seja
$\widehat{f}:\widehat{X}\to\widehat{Y}$ sua extensão às
completamentos. Então a aplicação $f\mapsto\widehat{f}$ é uma bijeção
\[
  \Hom_S(X,Y)
  \xrightarrow{\ \sim\ }
  \Hom_{\widehat{S}}(\widehat{X},\widehat{Y}).
\]
\end{theorem}

\begin{proof}
Demonstremos primeiro a injetividade. Suponhamos que $f,g:X\to Y$ são
$S$-morfismos tais que $\widehat{f}=\widehat{g}$. Existe então
uma vizinhança aberta $V$ de $X'=u^{-1}(S')$ sobre a qual $f$ e $g$
coincidem \sref[I]{I.10.9.4}. Como $u$ é fechado,
\sref{III.5.1.3.1} dá $V=X$, e portanto $f=g$.

Demonstremos agora a sobrejetividade. Seja
$h:\widehat{X}\to\widehat{Y}$ um
$\widehat{S}$-morfismo. Ponhamos
\[
  Z=X\times_S Y,
\]
e denotemos as projeções canônicas por $p:Z\to X$ e
$q:Z\to Y$. O pré-esquema $Z$ é de tipo finito sobre $S$
\sref[I]{I.6.3.4} e, portanto, noetheriano. Seja
$\widehat{Z}$ seu completamento ao longo de
\[
  Z'=p^{-1}\bigl(u^{-1}(S')\bigr).
\]
O pré-esquema formal $\widehat{Z}$ identifica-se canonicamente com
$\widehat{X}\times_{\widehat{S}}\widehat{Y}$, e suas projeções
sobre $\widehat{X}$ e $\widehat{Y}$ identificam-se com
$\widehat{p}$ e $\widehat{q}$ \sref[I]{I.10.9.7}.

Como $Y$ é separado sobre $S$, $\widehat{Y}$ é separado sobre
$\widehat{S}$ \sref[I]{I.10.15.7}. O morfismo grafo
\[
  \Gamma_h=(1_{\widehat{X}},h):
  \widehat{X}\longrightarrow\widehat{Z}
\]
é, por conseguinte, uma imersão fechada \sref[I]{I.10.15.4}. Seja
$\mathfrak{T}$ o subpré-esquema formal fechado de
$\widehat{Z}$ associado a esta imersão, e seja
$j:\mathfrak{T}\to\widehat{Z}$ a imersão canônica. Assim,
\[
  \Gamma_h=j\circ w,
\]
onde $w:\widehat{X}\to\mathfrak{T}$ é um isomorfismo
\sref[I]{I.10.14.3} cujo inverso é
$\widehat{p}\circ j$. Além disso, $\mathfrak{T}$ é próprio sobre
$\widehat{S}$ porque $\widehat{X}$ o é. Deduz-se de
\sref{III.5.1.8} que existe um subpré-esquema fechado $T$ de $Z$
tal que
\[
  \mathfrak{T}
  =\widehat{T}
  =T_{/(T\cap Z')},
\]
e que $j=\widehat{i}$, onde $i:T\to Z$ é a imersão
canônica \sref[I]{I.10.14.7}.

O morfismo $p\circ i:T\to X$ é um isomorfismo: sua extensão aos
completamentos é o isomorfismo
$\widehat{p}\circ\widehat{i}$, e basta aplicar
\sref{III.4.6.8}, observando como antes que $S$ é a única
vizinhança de $S'$ em $S$.
\oldpage[III]{157}
Seja $g:X\to T$ o inverso de $p\circ i$, e ponhamos
\[
  f=q\circ i\circ g.
\]
Então $f:X\to Y$ tem por grafo $\Gamma_f=i\circ g$. Como
$\widehat{g}$ é o inverso de
$\widehat{p\circ i}$, tem-se
\[
  \widehat{\Gamma_f}
  =\widehat{i}\circ\widehat{g}
  =j\circ w
  =\Gamma_h.
\]
Mas $\widehat{\Gamma_f}=\Gamma_{\widehat{f}}$
\sref[I]{I.10.9.8}; portanto, $h=\widehat{f}$, o que termina a
demonstração.
\end{proof}

Em linguagem categórica, o funtor
$X\mapsto\widehat{X}$ é, portanto, plenamente fiel
\sref[0]{0.8.1.6} da categoria dos pré-esquemas próprios sobre
$\Spec(A)$ para a categoria dos pré-esquemas formais próprios sobre
$\Spf(A)$, para todo anel noetheriano ádico $A$. Fornece, por
conseguinte, uma equivalência entre a primeira categoria e uma
subcategoria plena da segunda. Os objetos desta subcategoria se
denominarão doravante \emph{pré-esquemas formais algebrizáveis}.
Para tal pré-esquema formal $\mathfrak{X}$, existe um pré-esquema
ordinário $X$, próprio sobre $\Spec(A)$ e determinado salvo isomorfismo
único, tal que $\mathfrak{X}\simeq\widehat{X}$.

\begin{env}[5.4.2]
\label{III.5.4.2}
\emph{Escólio.}
Com a notação de \sref{III.5.4.1}, ponhamos
\[
  S_n=\Spec(A/\mathfrak{J}^{n+1}),\qquad
  X_n=X\times_S S_n,\qquad
  Y_n=Y\times_S S_n.
\]
Deduz-se de \sref{III.5.4.1} e \sref[I]{I.10.12.3} que
dar um $S$-morfismo $f:X\to Y$ equivale a dar um sistema indutivo
ádico $(S_n)$ \sref[I]{I.10.12.2} de
$S_n$-morfismos $f_n:X_n\to Y_n$.
\end{env}

\begin{remark}[5.4.3]
\label{III.5.4.3}
Contrariamente ao que poderia sugerir o teorema de existência
\sref{III.5.1.6}, existem pré-esquemas formais próprios sobre $\Spf(A)$
que \emph{não são algebrizáveis}, assim como existem espaços analíticos
compactos que não procedem de variedades algébricas complexas. Mais
adiante encontraremos tais pré-esquemas na «teoria de moduli»,
que, quando o corpo base é $\mathbb{C}$, trata precisamente das
variações infinitesimais da estrutura complexa de uma variedade
algébrica completa; é bem sabido que tais variações podem dar
lugar a variedades analíticas que não são algébricas.
\end{remark}

\begin{proposition}[5.4.4]
\label{III.5.4.4}
Seja $A$ um anel noetheriano ádico, ponhamos
$\mathfrak{S}=\Spf(A)$, e sejam
\[
  g:\mathfrak{X}\to\mathfrak{S},
  \qquad
  h:\mathfrak{Y}\to\mathfrak{S}
\]
dois morfismos próprios de pré-esquemas formais. Seja
$f:\mathfrak{X}\to\mathfrak{Y}$ um
$\mathfrak{S}$-morfismo. Se $f$ é finito e
$\mathfrak{Y}$ é algebrizável, então $\mathfrak{X}$ é
algebrizável.
\end{proposition}

\begin{proof}
As hipóteses sobre $g$ e $h$ já implicam que $f$ é próprio
\sref{III.3.4.1}; para que $f$ seja finito, basta que a fibra
$f^{-1}(y)$ seja finita para todo $y\in\mathfrak{Y}$
\sref[0]{0.4.8.11}. Por hipótese,
\[
  \sh{B}=f_*(\sh{O}_{\mathfrak{X}})
\]
é uma $\sh{O}_{\mathfrak{Y}}$-álgebra coerente
\sref[0]{0.4.8.6}. Escrevamos
$\mathfrak{Y}=\widehat{Y}$ e $h=\widehat{w}$, onde
$w:Y\to\Spec(A)$ é um morfismo próprio de pré-esquemas ordinários. O
teorema de existência fornece uma $\sh{O}_Y$-álgebra coerente $\sh{C}$
tal que $\sh{B}=\widehat{\sh{C}}$. Ponhamos
\[
  X=\Spec(\sh{C}),
\]
e seja $u:X\to Y$ o morfismo estrutural. A definição de
$\mathfrak{X}$ a partir de $\sh{B}$ \sref[0]{0.4.8.7} mostra então que
$\mathfrak{X}$ é canonicamente isomorfo a $\widehat{X}$ e
que $f$ se identifica com $\widehat{u}$; basta verificá-lo
quando $Y$ é afim.
\end{proof}

A Proposição \sref{III.5.1.8} é um caso particular de
\sref{III.5.4.4}.

\begin{theorem}[5.4.5]
\label{III.5.4.5}
Seja $A$ um anel noetheriano ádico, seja $\mathfrak{J}$ um ideal
de definição de $A$, ponhamos
\[
  S=\Spec(A),\qquad
  \mathfrak{S}=\widehat{S}=\Spf(A),
\]
e seja $f:\mathfrak{X}\to\mathfrak{S}$ um morfismo próprio de
pré-esquemas formais. Ponhamos
\[
  S_k=\Spec(A/\mathfrak{J}^{k+1}),\qquad
  X_k=\mathfrak{X}\times_{\mathfrak{S}}S_k,
\]
e, para todo $\sh{O}_{\mathfrak{X}}$-módulo $\sh{F}$, ponhamos
\[
  \sh{F}_k
  =\sh{F}\otimes_{\sh{O}_{\mathfrak{X}}}\sh{O}_{X_k}
  =\sh{F}/\mathfrak{J}^{k+1}\sh{F}.
\]
\oldpage[III]{158}
Seja $\sh{L}$ um
$\sh{O}_{\mathfrak{X}}$-módulo invertível, e suponhamos que
\[
  \sh{L}_0=\sh{L}/\mathfrak{J}\sh{L}
\]
é um $\sh{O}_{X_0}$-módulo amplo. Então $\mathfrak{X}$ é
algebrizável. Além disso, se $X$ é um $S$-pré-esquema próprio tal
que $\mathfrak{X}\simeq\widehat{X}$, existe um
$\sh{O}_X$-módulo amplo $\sh{M}$ tal que
$\sh{L}\simeq\widehat{\sh{M}}$; em particular, $X$ é projetivo
sobre $S$.
\end{theorem}

\begin{proof}
Apliquemos \sref{III.5.2.3} com
$\sh{F}=\sh{O}_{\mathfrak{X}}$. Existe um inteiro $n_0$ tal
que, para $n\geq n_0$, o homomorfismo canônico
\[
  \Gamma(\mathfrak{X},\sh{L}^{\otimes n})
  \longrightarrow
  \Gamma(X_0,\sh{L}_0^{\otimes n})
\]
é sobrejetivo. Escolhamos $n\geq n_0$ suficientemente grande para que
$\sh{L}_0^{\otimes n}$ seja muito amplo para $S_0$
\sref[II]{II.4.5.10}. Como
$f_0:X_0\to S_0$ é próprio,
$\Gamma(X_0,\sh{L}_0^{\otimes n})$ é um $A$-módulo de tipo finito
\sref{III.3.2.1}. Existe, por conseguinte, um
$A$-submódulo de tipo finito
\[
  E\subset\Gamma(\mathfrak{X},\sh{L}^{\otimes n})
\]
cuja imagem em $\Gamma(X_0,\sh{L}_0^{\otimes n})$ é todo o
módulo.

Para todo $k\geq0$, consideremos o homomorfismo
\[
  u_k:
  E/\mathfrak{J}^{k+1}E
  \longrightarrow
  \Gamma(X_k,\sh{L}_k^{\otimes n})
\]
induzido pela injeção canônica
$E\to\Gamma(\mathfrak{X},\sh{L}^{\otimes n})$. O módulo
$(f_k)_*(\sh{L}_k^{\otimes n})$ é quase-coerente, e
\[
  \Gamma\bigl(S_k,(f_k)_*(\sh{L}_k^{\otimes n})\bigr)
  =
  \Gamma(X_k,\sh{L}_k^{\otimes n}).
\]
Assim, $u_k$ define um homomorfismo
\[
  \widetilde{u}_k:
  \widetilde{E/\mathfrak{J}^{k+1}E}
  \longrightarrow
  (f_k)_*(\sh{L}_k^{\otimes n}),
\]
e, portanto, por adjunção, um homomorfismo
\[
  \widetilde{u}_k^{\,\sharp}:
  f_k^*\bigl(\widetilde{E/\mathfrak{J}^{k+1}E}\bigr)
  \longrightarrow
  \sh{L}_k^{\otimes n}.
\]
Ponhamos
\[
  \sh{G}_k=
  f_k^*\bigl(\widetilde{E/\mathfrak{J}^{k+1}E}\bigr).
\]
Tem-se
$\sh{G}_0=\sh{G}_k/\mathfrak{J}\sh{G}_k$
\sref[I]{I.9.1.5}, e
\[
  \widetilde{u}_0^{\,\sharp}:
  \sh{G}_0/\mathfrak{J}\sh{G}_0
  \longrightarrow
  \sh{L}_0^{\otimes n}/
  \mathfrak{J}\sh{L}_0^{\otimes n}
\]
obtém-se de $\widetilde{u}_k^{\,\sharp}$ passando aos
quocientes.

Pela definição de $E$, $\widetilde{u}_0^{\,\sharp}$ é o
homomorfismo canônico
\[
  \sigma:
  f_0^*\bigl((f_0)_*(\sh{L}_0^{\otimes n})\bigr)
  \longrightarrow
  \sh{L}_0^{\otimes n}.
\]
Como $\sh{L}_0^{\otimes n}$ é muito amplo,
$\widetilde{u}_0^{\,\sharp}$ é sobrejetivo
\sref[II]{II.4.4.3}. O lema de Nakayama mostra então que todo
$\widetilde{u}_k^{\,\sharp}$ é sobrejetivo.

Cada $\widetilde{u}_k^{\,\sharp}$ define, portanto,
\sref[II]{II.4.2.2} um $S_k$-morfismo
\[
  g_k:X_k\longrightarrow
  P_k=\mathbf{P}(E/\mathfrak{J}^{k+1}E).
\]
Para $h\leq k$,
\[
  P_h=P_k\times_{S_k}S_h
\]
por \sref[II]{II.4.1.3}. As relações entre os
$\widetilde{u}_k^{\,\sharp}$ e \sref[II]{II.4.2.10} mostram que
$(g_k)$ é um sistema indutivo ádico $(S_k)$
\sref[I]{I.10.12.2}. Os $g_k$ definem, por conseguinte, um
$\mathfrak{S}$-morfismo de pré-esquemas formais
\[
  g:\mathfrak{X}\longrightarrow\mathfrak{P},
\]
onde $\mathfrak{P}$ é o limite indutivo do sistema $(P_k)$,
ou, equivalentemente, o completamento $\widehat{P}$ de
$P=\mathbf{P}(E)$. O fato de que
$\sh{L}_0^{\otimes n}$ seja muito amplo implica que $g_0$ é uma imersão fechada
\sref[II]{II.4.4.3}. Portanto, $g$ é uma imersão fechada de
pré-esquemas formais \sref[0]{0.4.8.10}, e $\mathfrak{X}$ é algebrizável
por \sref{III.5.1.8}.

O teorema de existência \sref{III.5.1.6} mostra agora que $\sh{L}$ é
o completamento $\widehat{\sh{M}}$ de um
$\sh{O}_X$-módulo invertível $\sh{M}$. Além disso,
$\sh{L}^{\otimes n}$ é o completamento de
$\sh{M}^{\otimes n}$ \sref[I]{I.10.8.10}, e os homomorfismos
$\widetilde{u}_k^{\,\sharp}$ definem um único homomorfismo
\[
  v:f^*(\widetilde{E})
  \longrightarrow
  \sh{M}^{\otimes n}
\]
\sref{III.5.1.7}. Como os
$\widetilde{u}_k^{\,\sharp}$ são sobrejetivos,
$\widehat{v}$ é sobrejetivo \sref[I]{I.10.11.5}, e portanto também
o é $v$ \sref{III.5.1.3}. Seja $r:X\to P$ o morfismo definido
por $v$ \sref[II]{II.4.2.2}. Sua extensão aos completamentos é
$g$. Como $g$ é uma imersão fechada, também o é $r$, por
\sref{III.5.1.8} e \sref{III.5.4.1}. Deduz-se que
$\sh{M}^{\otimes n}$ é muito amplo \sref[II]{II.4.4.6} e, portanto,
que $\sh{M}$ é amplo \sref[II]{II.4.5.10}.
\end{proof}

\begin{remark}[5.4.6]
\label{III.5.4.6}
Seja $A$ um anel noetheriano ádico, ponhamos
$\mathfrak{S}=\Spf(A)$, e seja
$f:\mathfrak{X}\to\mathfrak{S}$ um morfismo próprio de pré-esquemas
formais. Seja $\sh{N}$ o ideal coerente de
$\sh{O}_{\mathfrak{X}}$ tal que, para todo subconjunto aberto
formal afim $U$ de $\mathfrak{X}$,
\[
  \Gamma(U,\sh{N})
\]
é o nilradical de
$\Gamma(U,\sh{O}_{\mathfrak{X}})$. A existência deste ideal
deduz-se facilmente de \sref[I]{I.10.10.2} e do fato de que
todo homomorfismo de anéis $B\to C$ envia o nilradical de $B$ em
o de $C$. Seja $\mathfrak{X}'$ o subpré-esquema formal fechado de
$\mathfrak{X}$ definido por $\sh{N}$
\sref[I]{I.10.14.2}. Seria interessante saber se a
algebraizabilidade de $\mathfrak{X}'$ implica a algebraizabilidade do
próprio $\mathfrak{X}$.

Sem dúvida, chegar-se-ia a uma solução deste problema se se soubesse
classificar, por exemplo mediante invariantes
\oldpage[III]{159}
de natureza cohomológica, as \emph{extensões} de um feixe estrutural
$\sh{O}_{\mathfrak{X}}$---para um pré-esquema ordinário ou
formal---por um ideal de quadrado nulo. Equivalentemente, estas são as
$\sh{O}_{\mathfrak{X}}$-álgebras $\sh{A}$ tais que
$\sh{O}_{\mathfrak{X}}\simeq\sh{A}/\sh{J}$, onde $\sh{J}$ é um
ideal de quadrado nulo de $\sh{A}$.
\end{remark}


\subsection{Uma decomposição de certos pré-esquemas}
\label{subsection:III.5.5}

\begin{proposition}[5.5.1]
\label{III.5.5.1}
Seja $A$ um anel noetheriano ádico, seja $\mathfrak{J}$ um ideal
de definição de $A$, e ponhamos $Y=\Spec(A)$. Seja $f:X\to Y$ um
morfismo separado de tipo finito, e ponhamos
\[
  Y_0=\Spec(A/\mathfrak{J}),\qquad
  X_0=X\times_Y Y_0=f^{-1}(Y_0).
\]
Seja $Z_0$ um subconjunto aberto de $X_0$ que é próprio sobre $Y_0$.
Então existe um subconjunto aberto e fechado $Z$ de $X$, próprio sobre
$Y$, tal que
\[
  Z\cap X_0=Z_0.
\]
\end{proposition}

\begin{proof}
Por hipótese, existe um subconjunto aberto $T$ de $X$ tal que
$T\cap X_0=Z_0$. Seja $\widehat{T}$ o completamento ao longo de $Z_0$
do pré-esquema induzido por $X$ sobre $T$. O suporte de
$\sh{O}_{\widehat{T}}$ é $Z_0$, que é próprio sobre $Y_0$;
portanto, $\widehat{T}$ é próprio sobre
$\widehat{Y}=\Spf(A)$ \sref{III.3.4.1}.

Deduz-se de \sref{III.5.1.8} que existe um subpré-esquema fechado
$Z$ de $T$, próprio sobre $Y$, tal que, se
$i:Z\to T$ é a imersão canônica, então
\[
  \widehat{i}:\widehat{Z}\longrightarrow\widehat{T}
\]
é um isomorfismo; aqui $\widehat{Z}$ é o completamento de $Z$
ao longo de $Z_0$. Por \sref{III.4.6.8}, existe uma vizinhança aberta
$V$ de $Z_0$ em $T$ tal que
$i^{-1}(V)\to V$ é um isomorfismo. Mas $i^{-1}(V)$ é uma
vizinhança de $Z_0$ em $Z$ e, portanto, é igual a $Z$ por
\sref{III.5.1.3.1}. Assim, $Z$ é aberto em $T$, e portanto em $X$.
Também é fechado em $X$, pois $Z$ é próprio sobre $Y$ e $X$ é
separado sobre $Y$. Isto demonstra a proposição.
\end{proof}

\begin{corollary}[5.5.2]
\label{III.5.5.2}
Se $X_0$ é próprio sobre $Y_0$, então $X$ é a união de dois
subconjuntos abertos disjuntos $Z$ e $Z'$ tais que $Z$ é próprio sobre
$Y$ e contém $X_0$. Além disso, todo subconjunto fechado $P$ de $X$
que é próprio sobre $Y$ está contido em $Z$.
\end{corollary}

\begin{proof}
Apliquemos \sref{III.5.5.1} com $Z_0=X_0$, e ponhamos $Z'=X-Z$. Para a
última afirmação, $P\cap Z'$ é fechado em $P$ e, portanto, próprio
sobre $Y$. Se não fosse vazio, $f(P\cap Z')$ seria um subconjunto fechado
não vazio de $Y$ e, portanto, intersectaria $Y_0$
\sref{III.5.1.3.1}. Isto contradiz
$Z'\cap X_0=\varnothing$ e demonstra que $P\subset Z$.
\end{proof}

\bigskip
\noindent\emph{(Continuará.)}

% La página impresa 504 está en blanco.
\clearpage
\thispagestyle{empty}
\null
\clearpage
\thispagestyle{plain}

\oldpage[III]{161}
\section*{BIBLIOGRAFIA (continuação)}

\begin{flushleft}
[27] P. Gabriel, \emph{Des catégories abéliennes}, tese, Paris (1961).

\medskip
[28] J. E. Roos, \emph{Sur les foncteurs dérivés de $\varprojlim$.
Applications}, C. R. Acad. Sci. Paris, vol. CCLII, 1961, pp. 3702--3704.

\medskip
[29] H. Cartan, \emph{Séminaire de l'École Normale Supérieure},
13.er anho (1960--1961), exposição n.º 11.
\end{flushleft}

\clearpage
\oldpage[III]{162}
\section*{ÍNDICE De Notações}

\begin{flushleft}
$\Set$: 0, 8.1.1.

$h_X(Y)$, $h_X(u)$ ($X$, $Y$ objetos de uma categoria $\C$, $u$ um morfismo
de $\C$): 0, 8.1.1.

$h_w(Y)$ ($Y$ um objeto e $w$ um morfismo de $\C$): 0, 8.1.2.

$Z_k(E_r^{pq})$, $B_k(E_r^{pq})$ ($k\geq r$): 0, 11.1.1.

$Z_\infty(E_2^{pq})$, $B_\infty(E_2^{pq})$: 0, 11.1.1.

$K^\bullet$, $K_\bullet$ (complexos): 0, 11.2.1.

$K^{\bullet\bullet}$, $K_{\bullet\bullet}$ (bicomplexos): 0, 11.2.1.

$B^i(K^\bullet)$, $Z^i(K^\bullet)$, $H^i(K^\bullet)$,
$H^\bullet(K^\bullet)$: 0, 11.2.1.

$B_i(K_\bullet)$, $Z_i(K_\bullet)$, $H_i(K_\bullet)$,
$H_\bullet(K_\bullet)$: 0, 11.2.1.

$T(K^\bullet)$, $T(K_\bullet)$ ($T$ um funtor): 0, 11.2.1.

$E(K^\bullet)$ ($K^\bullet$ um complexo filtrado): 0, 11.2.2.

$K^{i,\bullet}$, $K^{\bullet,j}$, $Z_{\mathrm{II}}^p(K^{i,\bullet})$,
$B_{\mathrm{II}}^p(K^{i,\bullet})$, $H_{\mathrm{II}}^p(K^{i,\bullet})$,
$Z_{\mathrm I}^p(K^{\bullet,j})$, $B_{\mathrm I}^p(K^{\bullet,j})$,
$H_{\mathrm I}^p(K^{\bullet,j})$: 0, 11.3.1.

$H_{\mathrm{II}}^p(K^{\bullet\bullet})$,
$Z_{\mathrm I}^q(H_{\mathrm{II}}^p(K^{\bullet\bullet}))$,
$B_{\mathrm I}^q(H_{\mathrm{II}}^p(K^{\bullet\bullet}))$,
$H_{\mathrm I}^q(H_{\mathrm{II}}^p(K^{\bullet\bullet}))$,
$H^n(K^{\bullet\bullet})$: 0, 11.3.1.

$F_{\mathrm I}^p(K^{\bullet\bullet})$,
$F_{\mathrm{II}}^p(K^{\bullet\bullet})$, ${}'E(K^{\bullet\bullet})$,
${}''E(K^{\bullet\bullet})$: 0, 11.3.2.

$K_{i,\bullet}$, $K_{\bullet,j}$, $H_p^{\mathrm{II}}(K_{i,\bullet})$,
$H_p^{\mathrm I}(K_{\bullet,j})$, $H_p^{\mathrm{II}}(K_{\bullet\bullet})$,
$H_p^{\mathrm I}(K_{\bullet\bullet})$,
$H_q^{\mathrm I}(H_p^{\mathrm{II}}(K_{\bullet\bullet}))$,
$H_q^{\mathrm{II}}(H_p^{\mathrm I}(K_{\bullet\bullet}))$,
$H_n(K_{\bullet\bullet})$: 0, 11.3.5.

$\RR^\bullet T(K^\bullet)$ ($T$ um funtor covariante aditivo): 0, 11.4.3.

$\RR^\bullet T(K^\bullet,K'^{\bullet})$ ($T$ um bifuntor covariante
aditivo): 0, 11.4.6.

$\LL_\bullet T(K_\bullet)$ ($T$ um funtor covariante aditivo): 0, 11.6.2.

$\LL_\bullet T(K_\bullet,K'_\bullet)$ ($T$ um bifuntor covariante
aditivo): 0, 11.6.6.

$\LL_\bullet T(K_{\bullet\bullet})$ ($T$ um funtor covariante aditivo):
0, 11.7.2.

$|\sigma|$, $\Sigma(A)$, $\mathrm C_\bullet(A)$, $\mathrm L_\bullet(A)$
($A$ um conjunto finito, $\sigma$ um simplexo de $A$): 0, 11.8.1.

$\mathrm C^\bullet(A,B;\sh{S})$, $\mathrm L^\bullet(A,B;\sh{S})$
($A$, $B$ conjuntos finitos, $\sh{S}$ um sistema de coeficientes): 0, 11.8.4.

$\mathrm P^\bullet(A,B;\sh{S})$ ($A$, $B$ conjuntos finitos, $\sh{S}$ um
sistema de coeficientes): 0, 11.8.6.

$\mathrm C^\bullet(A,B;\sh{S}^\bullet)$,
$\mathrm L^\bullet(A,B;\sh{S}^\bullet)$ ($A$, $B$ conjuntos finitos,
$\sh{S}^\bullet$ um complexo de sistemas de coeficientes): 0, 11.8.8.

$\mathrm P^\bullet(A,B;\sh{S}^\bullet)$ ($A$, $B$ conjuntos finitos,
$\sh{S}^\bullet$ um complexo de sistemas de coeficientes): 0, 11.8.10.

$\chi(M^\bullet)$ ($M^\bullet$ um complexo finito de módulos de comprimento finito):
0, 11.10.1.

$\check C^\bullet(\mathfrak U,\sh{F})$ ($\sh{F}$ um feixe de grupos abelianos
sobre $X$, $\mathfrak U$ um recobrimento aberto de $X$): 0, 12.1.2.

$\RR^p f_*(\sh{F})$ ($f:X\to Y$ um morfismo de espaços anelados,
$\sh{F}$ um feixe de $\sh{O}_X$-módulos): 0, 12.2.1.

$\sh{H}^p(f,\sh{K}^\bullet)$,
$\sh{H}^p_f(\sh{K}^\bullet)$, ${}'E(f,\sh{K}^\bullet)$,
${}''E(f,\sh{K}^\bullet)$ ($f:X\to Y$ um morfismo de espaços anelados,
$\sh{K}^\bullet$ um complexo de $\sh{O}_X$-módulos): 0, 12.4.1.

$\mathbf H^\bullet(X,\sh{K}^\bullet)$,
$\mathbf H^\bullet(V,\sh{K}^\bullet)$ ($X$ um espaço anelado,
$\sh{K}^\bullet$ um complexo de $\sh{O}_X$-módulos, $V$ aberto em $X$):
0, 12.4.2.

$\operatorname{gr}^\bullet(A)$ ($A$ um sistema projetivo que satisfaz (ML)):
0, 13.4.2.

$E(A)$ ($A$ um objeto filtrado): 0, 13.6.4.

$K_\bullet(\mathbf f)$, $i_{\mathbf f}$ ($\mathbf f$ uma família finita de
elementos de um anel): III, 1.1.1.

$K^\bullet(\mathbf f,M)$, $K_\bullet(\mathbf f,M)$ ($\mathbf f$ uma família
finita de elementos de um anel $A$, $M$ um $A$-módulo): III, 1.1.2.

$H^\bullet(\mathbf f,M)$, $H_\bullet(\mathbf f,M)$ ($\mathbf f$ uma família
finita de elementos de um anel $A$, $M$ um $A$-módulo): III, 1.1.3.

$C^\bullet((\mathbf f),M)$, $H^\bullet((\mathbf f),M)$ ($\mathbf f$ uma
família finita de elementos de um anel $A$, $M$ um $A$-módulo): III, 1.1.6.

$H^\bullet(U,\sh{F}^{(*)})$,
$H^\bullet(\mathfrak U,\sh{F}^{(*)})$ ($\sh{F}$ um
$\sh{O}_X$-módulo quase-coerente, $U$ aberto em $X$, $\mathfrak U$ um recobrimento aberto de $X$):
III, 2.1.1.
\end{flushleft}

\clearpage
\oldpage[III]{163}
\section*{ÍNDICE TERMINOLÓGICO}

\begin{flushleft}
Termo de convergência de uma sucessão espectral: 0, 11.1.1.

Algebrizável ($\sh{O}_X$-módulo): III, 5.2.1.

Algebrizável (esquema formal): III, 5.4.2.

Analiticamente íntegro (anel): III, 4.3.6.

Aplicação quase-compacta: 0, 9.1.1.

Aumento de uma resolução: 0, 11.4.1.

Característica de Euler--Poincaré de um complexo: 0, 11.10.1.

Característica de Euler--Poincaré de um $\sh{O}_X$-módulo coerente
($X$ projetivo sobre $\Spec(A)$, $A$ um anel artiniano): III, 2.5.1.

Cocadeia bialternada: 0, 11.8.4.

Complexo definido por um bicomplexo: 0, 11.3.1.

Complexo de álgebra exterior: III, 1.1.1.

Condição (ML): 0, 13.1.2.

Construtível (subconjunto, conjunto): 0, 9.1.2.

Construtível (função): 0, 9.3.1.

Produto cup: 0, 12.1.2.

Dihomomorfismo para estruturas algébricas sobre categorias: 0, 8.2.1.

Exato (subconjunto) em uma categoria abeliana: III, 3.1.1.

Essencialmente constante (sistema projetivo): 0, 13.4.2.

Fatoração de Stein: III, 4.3.3.

Codiscreta, cosseparada, discreta, exaustiva, finita, separada
(filtração): 0, 11.1.3.

Final (objeto) de uma categoria: 0, 8.1.10.

Finito (morfismo) de pré-esquemas formais: III, 4.8.2.

Finito (pré-esquema formal) sobre $\mathfrak Y$,
$\mathfrak Y$-finito (pré-esquema formal): III, 4.8.2.

Funtor covariante canônico $\C\to\CHom(\C\op,\Set)$: 0, 8.1.2.

Género aritmético: III, 2.5.1.

Geometricamente conexa (fibra): III, 4.3.4.

Homomorfismo de estruturas algébricas sobre categorias: 0, 8.2.1.

Hipercohomologia de um funtor com respeito a um complexo $K^\bullet$:
0, 11.4.3.

Hipercohomologia de um bifuntor com respeito a dois complexos
$K^\bullet$, $K'^{\bullet}$: 0, 11.4.6.

Hiper-homologia de um funtor com respeito a um complexo $K_\bullet$:
0, 11.6.2.

Hiper-homologia de um bifuntor com respeito a dois complexos $K_\bullet$,
$K'_\bullet$: 0, 11.6.6.

Hiper-homologia de um funtor com respeito a um bicomplexo: 0, 11.7.2.

Localmente construtível (conjunto): 0, 9.1.11.

Lei de composição externa (interna) sobre os objetos de uma categoria: 0, 8.2.1.

\oldpage[III]{164}

$S$-$\C$-módulo filtrado, $\operatorname{gr}^\bullet(S)$-$\C$-módulo graduado,
$\operatorname{gr}^{\bullet\bullet}(S)$-$\C$-módulo bigraduado: 0, 13.6.6.

Morfismo de sucessões espectrais: 0, 11.1.2.

Número geométrico de componentes conexas de uma fibra: III, 4.3.4.

Objeto grupo em $\C$, $\C$-grupo, $\C$-anel, $\C$-módulo: 0, 8.2.3.

Objeto de bordos, objeto de cobordes, objeto de cociclos, objeto de
ciclos: 0, 11.2.1.

Objeto de imagens universais: 0, 13.1.1.

Objeto graduado associado a um sistema projetivo que satisfaz (ML): 0, 13.4.2.

Objeto graduado limitado inferiormente (superiormente): 0, 11.2.1.

Subconjunto próprio (sobre $\mathfrak Y$) de um pré-esquema formal: III, 3.4.1.

Plena (subcategoria): 0, 8.1.5.

Plenamente fiel (funtor): 0, 8.1.5.

Polinômio de Hilbert: III, 2.5.3.

Próprio (morfismo) de pré-esquemas formais: III, 3.4.1.

Representável (funtor): 0, 8.1.8.

Cohomológica, direita, esquerda, homológica, injetiva, livre, plana, projetiva
(resolução): 0, 11.4.1.

Resolução de Cartan--Eilenberg direita, injetiva: 0, 11.4.2.

Resolução de Cartan--Eilenberg esquerda, projetiva: 0, 11.6.1.

Retrocompacto (conjunto): 0, 9.1.1.

Estrito (sistema projetivo): 0, 13.4.2.

Sucessão espectral em uma categoria abeliana: 0, 11.1.1.

Fracamente convergente, regular, coregular, birregular (sucessão espectral):
0, 11.1.3.

Sucessão espectral degenerada: 0, 11.1.6.

Sucessão espectral de um funtor relativa a um objeto filtrado: 0, 13.6.4.

Sucessões espectrais de um bicomplexo: 0, 11.3.2 e 11.3.5.

Sucessões espectrais de hipercohomologia de um funtor com respeito a um complexo:
0, 11.4.3.

Sucessões espectrais de hiper-homologia de um funtor com respeito a um complexo:
0, 11.6.2.

Sucessões espectrais de hiper-homologia de um funtor com respeito a um bicomplexo:
0, 11.7.2.

Sistema de coeficientes: 0, 11.8.4.

Unirrama (anel, ponto): III, 4.3.6.

Universalmente aberto (morfismo): III, 4.3.9.
\end{flushleft}

\clearpage
\oldpage[III]{165}
\section*{ÍNDICE GERAL}

\begin{flushleft}
\textsc{Capítulo 0.} --- \textbf{Preliminares (continuação)} \dotfill 5

\medskip
\S~8. Funtores representáveis \dotfill 5

8.1. Funtores representáveis \dotfill 5

8.2. Estruturas algébricas em categorias \dotfill 9

\medskip
\S~9. Conjuntos construtíveis \dotfill 12

9.1. Conjuntos construtíveis \dotfill 12

9.2. Subconjuntos construtíveis de espaços noetherianos \dotfill 14

9.3. Funções construtíveis \dotfill 16

\medskip
\S~10. Suplemento sobre os módulos planos \dotfill 17

10.1. Relações entre módulos planos e módulos livres \dotfill 17

10.2. Critérios locais de planicidade \dotfill 18

10.3. Existência de extensões planas de anéis locais \dotfill 20

\medskip
\S~11. Suplemento sobre álgebra homológica \dotfill 23

11.1. Lembrete sobre as sucessões espectrais \dotfill 23

11.2. A sucessão espectral de um complexo filtrado \dotfill 27

11.3. As sucessões espectrais de um bicomplexo \dotfill 29

11.4. Hipercohomologia de um funtor com respeito a um complexo $K^\bullet$
\dotfill 32

11.5. Passagem ao limite indutivo em hipercohomologia \dotfill 35

11.6. Hiper-homologia de um funtor com respeito a um complexo $K_\bullet$
\dotfill 39

11.7. Hiper-homologia de um funtor com respeito a um bicomplexo
$K_{\bullet\bullet}$ \dotfill 41

11.8. Suplemento sobre a cohomologia dos complexos simpliciais \dotfill 43

11.9. Um lema sobre complexos de tipo finito \dotfill 46

11.10. Característica de Euler--Poincaré de um complexo de módulos de comprimento finito
\dotfill 48

\medskip
\S~12. Suplemento sobre a cohomologia de feixes \dotfill 49

12.1. Cohomologia de feixes de módulos sobre espaços anelados \dotfill 49

12.2. Imagens diretas superiores \dotfill 57

12.3. Suplementos sobre os funtores Ext de feixes \dotfill 60

12.4. Hipercohomologia do funtor de imagem direta \dotfill 62

\oldpage[III]{166}

\S~13. Limites projetivos em álgebra homológica \dotfill 64

13.1. A condição de Mittag--Leffler \dotfill 64

13.2. A condição de Mittag--Leffler para grupos abelianos \dotfill 65

13.3. Aplicação: cohomologia de um limite projetivo de feixes \dotfill 68

13.4. Condição de Mittag--Leffler e objetos graduados associados a sistemas
projetivos \dotfill 69

13.5. Limites projetivos de sucessões espectrais de complexos filtrados
\dotfill 71

13.6. Sucessão espectral de um funtor relativa a um objeto dotado de uma
filtração finita \dotfill 73

13.7. Funtores derivados de um limite projetivo de argumentos \dotfill 75

\medskip
\textsc{Capítulo III.} --- \textbf{Estudo cohomológico dos feixes coerentes}
\dotfill 81

\medskip
\S~1. Cohomologia dos esquemas afins \dotfill 82

1.1. Lembrete sobre o complexo de álgebra exterior \dotfill 82

1.2. Cohomologia de \v Cech de um recobrimento aberto \dotfill 85

1.3. Cohomologia de um esquema afim \dotfill 88

1.4. Aplicação à cohomologia de pré-esquemas quaisquer \dotfill 89

\medskip
\S~2. Estudo cohomológico dos morfismos projetivos \dotfill 95

2.1. Cálculos explícitos de certos grupos de cohomologia \dotfill 95

2.2. O teorema fundamental dos morfismos projetivos \dotfill 100

2.3. Aplicação a feixes graduados de álgebras e módulos \dotfill 102

2.4. Uma generalização do teorema fundamental \dotfill 107

2.5. Característica de Euler--Poincaré e polinômio de Hilbert \dotfill 109

2.6. Aplicação: critérios de amplitude \dotfill 111

\medskip
\S~3. Teorema de finitude para os morfismos próprios \dotfill 115

3.1. O lema de dévissage \dotfill 115

3.2. O teorema de finitude: caso dos esquemas ordinários \dotfill 116

3.3. Generalização do teorema de finitude (esquemas ordinários) \dotfill 118

3.4. Teorema de finitude: caso dos esquemas formais \dotfill 119

\medskip
\S~4. O teorema fundamental dos morfismos próprios. Aplicações \dotfill 122

4.1. O teorema fundamental \dotfill 122

4.2. Casos particulares e variantes \dotfill 129

4.3. Teorema de conexão de Zariski \dotfill 130

4.4. «Teorema principal» de Zariski \dotfill 135

4.5. Completamentos de módulos de homomorfismos \dotfill 138

4.6. Relações entre morfismos formais e morfismos ordinários \dotfill 139

4.7. Um critério de amplitude \dotfill 145

4.8. Morfismos finitos de pré-esquemas formais \dotfill 146

\oldpage[III]{167}

\S~5. Um teorema de existência para os feixes algébricos coerentes \dotfill 149

5.1. Enunciado do teorema \dotfill 149

5.2. Demonstração do teorema de existência: caso projetivo e quase-projetivo
\dotfill 151

5.3. Demonstração do teorema de existência: caso geral \dotfill 154

5.4. Aplicação: comparação entre morfismos de pré-esquemas ordinários e
morfismos de pré-esquemas formais. Pré-esquemas formais algebrizáveis \dotfill 156

5.5. Uma decomposição de certos pré-esquemas \dotfill 159

\medskip
\textsc{Bibliografia (continuação)} \dotfill 161

\textsc{Índice de notações} \dotfill 162

\textsc{Índice terminológico} \dotfill 163

\begin{flushright}
\emph{Recebido em 28 de junho de 1961.}
\end{flushright}
\end{flushleft}
